\chapter{Prenatal Panel}

\section*{What Is This Test?}
The prenatal panel, also known as the first-trimester screen, second-trimester maternal serum screen (quadruple screen), or sequential/integrated screen, is a set of blood tests and ultrasound measurements performed during pregnancy to screen for fetal chromosomal aneuploidy (trisomy 21, 18) and open neural tube defects \cite{acog2020screening}. The first-trimester screen (performed at 11--13 weeks 6 days) combines maternal serum markers (PAPP-A --- pregnancy-associated plasma protein A, and free beta-hCG or total hCG) with the ultrasound measurement of the fetal nuchal translucency (NT). The second-trimester quadruple screen (performed at 15--22 weeks) measures maternal serum alpha-fetoprotein (MSAFP), hCG (or free beta-hCG), unconjugated estriol (uE3), and inhibin A. The integrated screen combines the first-trimester and second-trimester results into a single risk estimate (disclosed after the second-trimester draw). The sequential screen discloses the first-trimester result, and if high risk, offers CVS; if low risk, the second-trimester draw is performed. Biochemical markers are measured by automated chemiluminescent immunoassay on platforms including PerkinElmer Delfia Xpress (AutoDELFIA), Beckman Coulter Access, Siemens Immulite, or Roche Cobas. Results are reported as multiples of the median (MoM) adjusted for gestational age, maternal weight, race, smoking status, insulin-dependent diabetes, and number of fetuses. The combined first-trimester screen has a detection rate of 82--87\% for trisomy 21 at a 5\% false-positive rate \cite{malone2005first}. The integrated screen has a detection rate of 94--96\% at a 5\% false-positive rate \cite{wald2003suruss}.

\section*{Alternative Names}
First-Trimester Screen; Quadruple Screen (Quad Screen); Triple Screen; Integrated Screen; Sequential Screen; Maternal Serum Screen; MSAFP Screen; Prenatal Aneuploidy Screen; Prenatal Genetic Screening

\section*{Principle}
Maternal serum biochemical markers are measured by time-resolved fluoroimmunoassay (DELFIA) or chemiluminescent immunoassay: pregnancy-associated plasma protein A (PAPP-A) by sandwich immunoassay; free beta-hCG by sandwich immunoassay; MSAFP by sandwich immunoassay; unconjugated estriol (uE3) by competitive immunoassay; and inhibin A by sandwich immunoassay. Nuchal translucency is measured by certified sonographers using standardized technique: a mid-sagittal view of the fetus with the fetus in a neutral position, measuring the maximum thickness of the subcutaneous translucent space at the back of the fetal neck (NT) with calipers, NT measured to the nearest 0.1 mm. Results are entered into a risk calculation algorithm (e.g., PerkinElmer LifeCycle, Fetal Medicine Foundation [FMF] algorithm) that combines maternal age, biochemical markers in MoM, and NT measurement to calculate the individualized risk.

\section*{History}
Elevated MSAFP was linked to fetal neural tube defects in 1972 by Brock and Sutcliffe \cite{brock1972alpha}. In 1984, Merkatz et al. demonstrated that low MSAFP was associated with trisomy 21. The triple screen (MSAFP + hCG + uE3) was introduced in the late 1980s--early 1990s. Nuchal translucency measurement as a marker for trisomy 21 was described by Nicolaides et al. in 1992 \cite{nicolaides1992fetal}. The first-trimester screen combining PAPP-A, free beta-hCG, and NT was validated in large multicenter studies (SURUSS, FASTER) in the 1990s--2000s \cite{nicolaides2005multicenter}. Inhibin A as a fourth second-trimester marker (quadruple screen) was added in the 1990s. The introduction of NIPT/cfDNA from 2011 onward has substantially changed prenatal screening algorithms, with many centers now offering NIPT as the primary screening test. However, traditional serum screening remains important, particularly for women who do not have access to NIPT, for detection of neural tube defects (not detected by NIPT), and for early detection of placental dysfunction (low PAPP-A).

\section*{Why This Test Is Ordered}
Screening for fetal trisomy 21 (Down syndrome) and trisomy 18 (Edwards syndrome) \cite{acog2016screening}. Screening for open neural tube defects (anencephaly, open spina bifida) by MSAFP measurement in the second trimester. Detection of pregnancies at risk for adverse outcomes related to placental insufficiency: low PAPP-A (<0.4 MoM in the first trimester) is associated with an increased risk of preeclampsia, fetal growth restriction, preterm birth, and stillbirth.

\section*{Primary vs Secondary Test}
This is a screening test (NOT diagnostic). A positive screening result must be followed by diagnostic testing (CVS or amniocentesis) for confirmation.

\section*{How This Test Is Performed}
First-trimester screen: maternal blood (SSR tube, gold-top) drawn at 11--13 weeks 6 days, combined with nuchal translucency ultrasound. Quadruple screen: maternal blood drawn at 15--22 weeks (optimal 16--18 weeks).

\section*{What Preparation Is Needed}
No special preparation is required, and no fasting is necessary. Correct gestational dating is essential: the first-trimester screen requires a maternal blood draw at 11 weeks through 13 weeks 6 days of gestation performed on the same day as the nuchal translucency ultrasound, and the second-trimester quadruple screen is drawn at 15--22 weeks (optimally 16--18 weeks). The patient should provide accurate information on menstrual dates and ultrasound dating, and should report insulin-dependent diabetes and smoking status, which are used in risk adjustment. No bladder filling or special diet is needed for the blood draw.

\section*{Contraindications and Precautions}
There are no absolute contraindications to the blood draw. The screen is not performed beyond the gestational-age window for each component, and women who have already undergone diagnostic testing do not need screening. Precautions: This is a screening test, not a diagnostic test, and a screen-positive result must be confirmed by chorionic villus sampling (CVS) or amniocentesis. Interpretation cautions: (1) Twin and higher-order multiple gestations require specialized risk algorithms, and marker concentrations differ from singletons. (2) Pregnancies conceived by in vitro fertilization or with oocyte donation, and pregnancies in women of different racial or ethnic groups, with insulin-dependent diabetes, or with vaginal bleeding require adjustments to the multiples-of-the-median (MoM) calculations. (3) Nuchal translucency measurement is operator-dependent and must be performed by certified sonographers using a standardized technique. (4) An elevated MSAFP may reflect gestational dating errors, multiple gestation, abdominal wall defects, or fetal demise in addition to open neural tube defects, and requires targeted ultrasound evaluation.

\section*{Risks}
Minimal risk. The test requires only a standard maternal venipuncture and a transabdominal ultrasound and carries no risk of miscarriage or injury to the fetus, unlike invasive diagnostic procedures (chorionic villus sampling or amniocentesis). Slight pain, bruising, or, rarely, hematoma at the venipuncture site.

\section*{What Happens After the Test}
Results are typically reported within 1--2 weeks as an individualized risk estimate (for example, 1 in 500) derived from the MoM values and gestational age. A low-risk result allows routine prenatal care to continue, including the 18--22 week anatomy ultrasound. A screen-positive result (risk above the laboratory cutoff, nuchal translucency $\geq$3.5 mm, or MSAFP >2.5 MoM) triggers genetic counseling and discussion of diagnostic testing (CVS at 10--13 weeks or amniocentesis at 15--20 weeks) or second-tier cfDNA (NIPT) screening. Women with a low PAPP-A (<0.4 MoM) receive increased surveillance for preeclampsia and fetal growth restriction even if the aneuploidy screen is negative.

\section*{Understanding Results}

\subsection*{Normal/Low-Risk Screen}
The calculated risk of trisomy 21 is less than a predetermined cutoff (e.g., <1:300 or <1:270, laboratory-dependent). The calculated risk of trisomy 18 is low. MSAFP is 0.5--2.5 MoM (no evidence of neural tube defect). Routine prenatal care continues.

\subsection*{High-Risk First-Trimester Screen}
Increased risk for trisomy 21 ($\geq$1:300 cutoff). NT $\geq$3.5 mm ($\geq$99th percentile for gestational age). Low PAPP-A (<0.4 MoM) with elevated free beta-hCG. Low PAPP-A AND low free beta-hCG suggests increased risk for trisomy 18 and trisomy 13. Referral for genetic counseling and discussion of diagnostic testing (CVS at 10--13 weeks), or cfDNA (NIPT) as a second-tier screening test before proceeding to invasive testing.

\subsection*{High-Risk Quadruple Screen}
Trisomy 21 pattern: Low MSAFP (<0.5 MoM), elevated hCG (>2.0 MoM), low uE3 (<0.5 MoM), and elevated inhibin A (>2.0 MoM). Trisomy 18 pattern: Low MSAFP, low hCG, low uE3 (all three low). Neural tube defect: Elevated MSAFP >2.5 MoM with a normal or unchanged hCG, uE3, inhibin A. Confirmation by detailed fetal anatomy ultrasound at 18--22 weeks. Referral for amniocentesis (15--20 weeks).

\subsection*{Low PAPP-A (<0.4 MoM) --- Increased Risk of Placental Dysfunction}
Increased monitoring for fetal growth restriction (serial growth ultrasounds), preeclampsia (blood pressure monitoring, urine protein), and preterm birth. Low-dose aspirin 81--162 mg/day initiated before 16 weeks is recommended to reduce the risk of preeclampsia in women with risk factors.

\section*{Normal Results}
First trimester (11--13.6 weeks): PAPP-A 0.5--2.0 MoM, free beta-hCG 0.5--2.0 MoM, NT <3.0 mm (varies by gestational age and crown-rump length). Second trimester (15--22 weeks): MSAFP 0.5--2.5 MoM, hCG 0.5--2.5 MoM, uE3 0.5--2.5 MoM, inhibin A 0.5--2.5 MoM. Trisomy 21 screen risk <1:300 (low risk, screen negative). Trisomy 18 screen risk low.

\section*{Follow-up Tests}
cfDNA (NIPT) as second-tier screening for women with a high-risk biochemical screen before proceeding to invasive testing. Chorionic villus sampling (CVS) with karyotype or CMA at 10--13 weeks for high-risk first-trimester screen. Amniocentesis with karyotype or CMA at 15--20 weeks for high-risk second-trimester screen. Fetal anatomy ultrasound (detailed) at 18--22 weeks. Growth ultrasounds every 3--4 weeks starting at 24--28 weeks if low PAPP-A. MSAFP amniotic fluid level and acetylcholinesterase if elevated MSAFP and normal anatomy ultrasound (for rare neural tube defects not visible on ultrasound).

\section*{References}
\bibliographystyle{unsrtnat}
\bibliography{references}

\clearpage
